% Source-derived chapter 04: Applications to Cremona groups
% Original source: motivic-invariants-birational-maps
\chapter{Applications to Cremona groups}
\label{motivic-invariants-birational-maps-chapter-04}
\source{motivic-invariants-birational-maps}

\begin{remark}[Source section]
\label{motivic-invariants-birational-maps-chapter-04-source}
\source{motivic-invariants-birational-maps}
\source{motivic-invariants-birational-maps}
This chapter reproduces the corresponding section of the retrieved
LaTeX source, including its definitions, statements, examples, and
proofs. It has not yet been translated into Lean.
\notready
\end{remark}

% The source section title was: Applications to Cremona groups
\label{sec:cremona}
\source{motivic-invariants-birational-maps}

Let $X/\k$ be a variety.

\begin{definition}
\source{motivic-invariants-birational-maps}
\notready\label{Def-reg}
\source{motivic-invariants-birational-maps}
A birational automorphism $\phi \in \Bir(X)$ is pseudo-regularizable
if $\phi = \alpha^{-1} \circ \gamma \circ \alpha$ where
\begin{itemize}
    \item $\alpha: X \dashrightarrow X'$ is a birational map
    (we do not assume $X'$ is smooth nor projective);
    \item $\gamma \in \Bir(X')$ is a pseudo-automorphism, namely
    an isomorphism in codimension one.
\end{itemize}

If moreover the map $\gamma$
is biregular, we call $\phi$ a regularizable map.
\end{definition}

Pseudo-regularizable maps have been
studied in \cite{Cantat-Cornulier, Lonjou-Urech}
in the context of the actions of Cremona
groups on various combinatorial complexes.

\begin{example}
\source{motivic-invariants-birational-maps}
\notready \label{ex-regbir} \hfill
\source{motivic-invariants-birational-maps}
\begin{enumerate}
    \item Regular automorphisms, or more generally
birational maps
$\phi \in \Bir(X)$ which have an invariant dense open subset
$U \subset X$, are regularizable.
\item Finite order elements of $\Bir(X)$ are regularizable.
Indeed, given such 
an element $\phi \in \Bir(X)$, the complement 
\[
U := X \setminus \bigcup_{k = 1}^{\ord(\phi)-1} \Ex(\phi^k),
\]
is an invariant dense open subset, and 
we apply (1) to conclude.
By Weil's regularization theorem (see also~\cite{Kraft}), 
they are even projectively regularizable,
namely there exists a regularization $\alpha : X \dto X'$
with projective $X'$.
\end{enumerate}
\end{example}


Let us write 
$\Bir(X)^{\psreg}$ for the normal subgroup
generated by pseudo-regularizable elements.

\begin{lemma}
\source{motivic-invariants-birational-maps}
\notready\label{lem:regularizable}
\source{motivic-invariants-birational-maps}
We have 
\[
\Bir(X)^{\psreg} \subset \Ker(c|_{\Bir(X)}).
\]
\end{lemma}
\begin{proof}
It suffices to show that
for any pseudo-regularizable element
$\phi \in \Bir(X)^{\psreg}$ we have
$c(\phi) = 0$.
By definition, for any such element 
we have $\phi = \alpha^{-1} \circ \gamma \circ \alpha$ where $\gamma$ is a pseudo-automorphism.
Thus $c(\gamma) = 0$ and 
\begin{equation*}
c(\phi) = -{c}(\alpha) + {c}(\gamma) + {c}(\alpha) = 0.
\end{equation*}
\end{proof}


\begin{theorem}
\source{motivic-invariants-birational-maps}
\notready\label{thm:bir-generation}
\source{motivic-invariants-birational-maps}
Assume that $X/\k$ is a variety
that
satisfies one of the following conditions:
\begin{enumerate}
    \item[(a)] $\k$ is a number field,
    or a function field over a number field,
    over a finite field, or over an algebraically closed field,
    and $X$ is birational to $\P^3 \times W$ for a
    separably rationally
    connected variety $W$
    (e.g. $X = \P^n$, $n \ge 3$).
    \item[(b)] $\k \subset \C$ is any subfield
    and $X$ is birational to $\P^4 \times W$ for a  rationally
    connected variety $W$
    (e.g. $X = \P^n$, $n \ge 4$).
    \item[(c)] $\k$ is any infinite
    field
    and $X$ is birational to $\P^5 \times W$ for a separably rationally
    connected variety $W$ 
    (e.g. $X = \P^n$, $n \ge 5$).
\end{enumerate}
We have $c|_{\Bir(X)} \ne 0$. 
As a consequence, 
$\Bir(X)$ is not generated by pseudo-regularizable elements. 
\end{theorem}

In particular, $\Bir(X)$ is not generated by 
any collection of elements from Example~\ref{ex-regbir}.



\begin{proof}
By Lemma~\ref{lem:regularizable}, 
it suffices
to explain that $c|_{\Bir(X)}$ is nonzero
for $\k$ and $X$ as in
the statement of the theorem.

Let us first assume that $W = \Spec(\k)$.
In case (a) (resp. (b) and (c)) we 
use the motivically non-trivial L-link constructed in Proposition~\ref{prop:elliptic} 
(resp. Theorem~\ref{thm:K3} and Theorem~\ref{thm:CY3}) ,
which together with 
Lemma~\ref{lemma:L-link} gives a map $\phi \in \Bir(\P^3)$
(resp. $\Bir(\P^4)$, $\Bir(\P^5)$)
such that 
$$c(\phi) = [Z \times \P^1] - [Z' \times \P^1] \ne 0 \in \Z[\Bir/\k].$$
where $Z$ and $Z'$ are non-isomorphic
K-trivial
varieties of Picard number $1$.
For an arbitrary $W$ as in Theorem~\ref{thm:bir-generation},
it then follows from Corollary~\ref{cor-CYnSB}
that
$$c(\phi \times \Id_W) = [Z \times \P^1 \times W] - [Z' \times \P^1 \times W] 
\ne 0 \in \Z[\Bir/\k].$$
This proves the non-vanishing of $c|_{\Bir(X)}$, and Theorem~\ref{thm:bir-generation} follows.
\end{proof}




\begin{corollary}
\source{motivic-invariants-birational-maps}
\notready\label{cor:homZ}
\source{motivic-invariants-birational-maps}
Let $X/\k$ be as in cases (a) or (b) in Theorem \ref{thm:bir-generation}.
Then we have a surjective homomorphism
$
\Bir(X) \to A
$
where $A = \bigoplus_J \Z$
is a free abelian group
with $J$ a set 
of the same cardinality as $\k$.
In particular, $A$ is a direct summand
of $\Bir(X)^{\ab}$.
\end{corollary}


\begin{proof}
We use the same
L-links which appear in the proof of
Theorem~\ref{thm:bir-generation}.
By Theorem~\ref{thm:ell-curves}
in case (a), and Theorem~\ref{thm:K3}
in case (b),
these L-links are
parameterized by a set $I$ having the cardinality of $\k$.
 
Let $X_i$ and $Y_i$ be the centers of these links.
Since the $X_i$'s are mutually non-isomorphic $K$-trivial varieties of Picard rank one, their birational types are all distinct (see e.g.
Theorem~\ref{thm:BurtCY}).
Furthermore, $X_i$ and $Y_i$
are in a symmetric relation with each other.
Namely, in case (a), the relation
is $Y_i \simeq \Jac^2(X_i)$,
so that $X_i \simeq \Jac^2(Y_i)$,
because we work with index $5$ genus one curves, see \cite[2.2]{ShinderZhang}.
In case (b), the relation is:
$Y_{i, \C}$ is the 
unique nontrivial
Fourier--Mukai partner of $X_{i,\C}$ \cite[3.1]{HassettLai}.
Thus in either case,
$$J \cnec \Set{ \Set{X_i,Y_i} \subset I | i \in I}$$
forms a partition of $I$
and has the same cardinality as $\k$.


We let
$$A \cnec \bigoplus_{\Set{X_i,Y_i} \in J}\Z\left([X_i \times \P^1 \times W] - [Y_i \times \P^1 \times W]\right) \subset
\Z[\Bir/\k].$$
By construction this a direct summand 
of $\Z[\Bir/\k]$, and let $\pi: \Z[\Bir/\k] \to A$ be the  projection homomorphism.
For each $i$ there exists
$\phi \in \Bir(X)$ such that $$\pi(c(\phi)) = 
[X_i \times \P^1 \times W] - [Y_i \times \P^1 \times W].$$
Thus the composition
$\pi \circ c: \Bir(X) \to A$ is surjective.

The statement about the abelianization follows
because $\Bir(X)^{\ab}$ admits a surjective homomorphism onto the free $\Z$-module $A$.
This homomorphism admits a splitting, hence it is 
a projection onto a direct summand.
\end{proof}


\begin{remark}
\source{motivic-invariants-birational-maps}
\notready
(i) After the results of this paper were announced, Blanc, Schneider and Yasinsky have proved that for all $n \ge 4$, $\Cr_n(\C)$ 
admits a surjective homomorphism onto an uncountable free product of $\Z$
\cite[Theorem B]{BSY} (which is
a stronger statement
than what
the Corollary \ref{cor:homZ} says
for $X = \P^n$ and $\k = \C$), with a completely different method.

(ii) For $n = 3$ we do not know
if there exist nontrivial homomorphisms $\Cr_3(\C) \to \Z$
(or $\Cr_3(\ol{\Q}) \to \Z$).
Whether $\Cr_3(\C)$ is generated by involutions, or regularizable elements, is an 
outstanding open
question about Cremona groups.
\end{remark}




\appendix
